% Figure: a two-leg sling on a load whose centre of gravity is offset from
% the geometric midpoint. The offset throws unequal tension into the two
% legs (non-equalising sling); the FBD of the load shows the moment balance
% that sets the split. Numbers match the case study.
\begin{figure}[htbp]
\centering
\begin{tikzpicture}[
  leg/.style={draw=engblue, line width=1.4pt},
  ten/.style={draw=engred, -{Stealth}, very thick},
  wt/.style={draw=engblack, -{Stealth}, very thick},
  scale=1.0,
]
% ===== Left: the offset-CG scene =====
\begin{scope}[xshift=0cm]
  \coordinate (H) at (2.3,4.0);
  \draw[draw=engblack, very thick] (H) -- ++(0,0.5);
  \draw[draw=engblack, fill=white, thick] (2.3,4.65) circle (0.18);
  \coordinate (L) at (0.6,1.4);
  \coordinate (R) at (4.0,1.4);
  \draw[leg] (H) -- (L);
  \draw[leg] (H) -- (R);
  \draw[draw=engblack, fill=engblue!10, thick] (0.4,0.4) rectangle (4.2,1.4);
  \fill[engblack] (L) circle (0.06);
  \fill[engblack] (R) circle (0.06);
  \node[font=\scriptsize, anchor=south] at (L) {$L$};
  \node[font=\scriptsize, anchor=south] at (R) {$R$};
  % centre of gravity, offset toward R
  \coordinate (G) at (2.8,0.9);
  \fill[engred] (G) circle (0.07);
  \node[font=\scriptsize, anchor=north, engred] at (G) {$G$};
  \draw[wt] (G) -- ++(0,-1.0) node[anchor=north, font=\scriptsize] {$W$};
  % geometric midpoint marker
  \draw[draw=enggray, thin, dashed] (2.3,1.4) -- (2.3,-0.1);
  \node[font=\scriptsize, anchor=north, enggray] at (2.3,-0.1) {midpoint};
  \node[font=\small, anchor=north] at (2.3,-0.7) {(a) offset centre of gravity};
\end{scope}
% ===== Right: FBD of the load (moment balance) =====
\begin{scope}[xshift=7.0cm]
  \coordinate (L2) at (0,1.4);
  \coordinate (R2) at (3.4,1.4);
  \draw[draw=engblack, fill=engblue!10, thick] (-0.2,0.4) rectangle (3.6,1.4);
  \fill[engblack] (L2) circle (0.06);
  \fill[engblack] (R2) circle (0.06);
  \node[font=\scriptsize, anchor=north east] at (L2) {$L$};
  \node[font=\scriptsize, anchor=north west] at (R2) {$R$};
  % leg tensions (vertical components shown for the moment balance)
  \draw[ten] (L2) -- ++(-0.6,1.5) node[anchor=south, font=\scriptsize, engred] {$T_L$};
  \draw[ten] (R2) -- ++(0.6,1.5) node[anchor=south, font=\scriptsize, engred] {$T_R$};
  % weight at offset G
  \coordinate (G2) at (2.2,0.9);
  \fill[engred] (G2) circle (0.06);
  \draw[wt] (G2) -- ++(0,-1.0) node[anchor=north, font=\scriptsize] {$W$};
  % distance markers
  \draw[draw=enggray, thin, <->] (0,-0.1) -- (2.2,-0.1);
  \node[font=\scriptsize, anchor=north, enggray] at (1.1,-0.1) {$a$};
  \draw[draw=enggray, thin, <->] (2.2,-0.55) -- (3.4,-0.55);
  \node[font=\scriptsize, anchor=north, enggray] at (2.8,-0.55) {$b$};
  \node[font=\small, anchor=north] at (1.6,-1.0) {(b) FBD of the load};
\end{scope}
\end{tikzpicture}
\caption[A two-leg sling on a load with an offset centre of gravity]{(a) A
crate slung from two legs, with its centre of gravity $G$ displaced toward
the right attachment. (b) The free-body diagram of the load: each leg
applies its tension at its own attachment, and the weight $W$ acts at $G$, a
distance $a$ from the left attachment and $b$ from the right. Moments about
the right attachment give the left leg's vertical share, $T_{L,v} = Wb/(a+b)$,
and about the left attachment the right leg's, $T_{R,v} = Wa/(a+b)$. The leg
nearer the centre of gravity takes the larger share. A non-equalising sling
locks each leg to its computed length and lets the split stand; an
equalising sling (a single line reeved through a ring at the hook) forces
$T_L = T_R$ and instead tilts the load until the geometry restores moment
balance. The choice governs whether the load hangs level or hangs safe.}
\label{fig:vol03ch01:sling-cg-offset}
\end{figure}
